
\section{Related Work}
\label{related_work}
\textbf{Video generation.}
The mainstream video generation methods can be categorized into GAN-based~\cite{goodfellow2014generative, vgan, mocogan}, Diffusion-based~\cite{ho2020ddpm, zhou2022magicvideo,ho2022video, ho2022imagen,SVD, make-a-video,girdhar2023emu, li2023videogen, zhang2023show1,openai2024sora} and language-model-based~\cite{vaswani2017attention,kondratyuk2023videopoet,villegas2022phenaki,yu2023language}. Among them, Diffusion-based methods have recently gained the most popularity. Most Diffusion-based methods encode videos into latent space~\cite{rombach2022high} for efficient training and utilize progressive inference strategies~\cite{ho2022imagen, he2022lvdm, AlignYourLatents} to generate videos with high spatial-temporal resolution. With a new scalable Diffusion Transformer~\cite{peebles2023scalable} architecture, Sora~\cite{openai2024sora} has further pushed video generation to a new stage. Different from diffusion-based video generation methods, our work aims to explore and unleash the potentiality of language models for long video generation, as their ability for modeling long sequence and scaling up have been proved in NLP.


\textbf{Image and video generation with language models.}
Language models have recently been explored for visual generation, with most works focusing on tokenizing visual data into a form that can be processed by these models. Quantization techniques like VQ-VAE \cite{vqvae,esser2020taming} are commonly used, and transformers are employed to model the resulting tokens.
For image generation, autoregressive or masked transformers are prevalent \cite{ramesh2021dalle, yu2022scaling, chang2023muse, chang2022maskgit, yu2024spae, sun2024autoregressive, tian2024visual}. In short video generation, image-level or video-level tokenizers are utilized, incorporating spatial-temporal compression and causal structures. Transformers model the spatial-temporal relationships, with various techniques proposed, such as sparse attention, spatial-temporal attention, large-scale pre-training, and improved tokenization \cite{yan2021videogpt,ge2022long, wu2021godiva, hong2022cogvideo, yu2023magvit}. VideoPoet \cite{kondratyuk2023videopoet} stands out as a multimodal model using bidirectional attention for conditioning, while our method aligns better with the language model paradigm by using unidirectional attention for both text and video.
However, these short video generation models focus on producing 1-5 second clips, limiting their ability to capture complex events and maintain consistency over longer durations.


\textbf{Long video generation.}
Previous works have explored long video generation using various approaches. LongVideoGAN~\cite{brooks2022generating}, NUWA-XL~\cite{yin2023nuwa}, and GAIA-1~\cite{hu2023gaia} utilized GAN-based methods, diffusion-over-diffusion techniques, or world models but were limited to specific domains. More recently, video diffusion models have been extended for longer video generation. FreeNoise~\cite{qiu2023freenoise} and Gen-L~\cite{wang2023gen_L} focus on sampling noise vectors and aggregating overlapping short video segments, respectively, while StreamingT2V~\cite{streamingt2v} proposes an autoregressive approach with memory blocks for consistency and appearance preservation. In the language model domain, Phenaki~\cite{villegas2022phenaki} generates variable-length videos using a masked video transformer. Despite these advancements, generating long videos with rich motion dynamics, consistent appearance, and high visual quality in the open domain remains a challenge.